\section{Gallego QLoRA}

\subsection{Tarea: CALIDAD\_LENGUA}\label{tarea-calidad_lengua}
\begin{longtable}[]{@{}lll@{}}
\toprule\noalign{}
Métrica & Base & Entrenado \\
\midrule\noalign{}
\endhead
\bottomrule\noalign{}
\caption{Métricas de rendimiento para la tarea Calidad Lengua (Gallego QLoRA).}
\label{tab:metricas_calidad_lengua_Gallego QLoRA}
\endlastfoot
ttr & 0.48 & \textbf{\underline{0.80}} \\
entropy & \textbf{\underline{6.72}} & 4.53 \\
ngram\_overlap & 0.00 & 0.00 \\
freq\_target & 0.85 & \textbf{\underline{0.99}} \\
freq\_comparison\_es & 0.68 & 0.68 \\
freq\_comparison\_fr & 0.33 & \textbf{\underline{0.35}} \\
calidad & 0.29 & \textbf{\underline{0.45}} \\
\end{longtable}

\subsubsection{Ejemplo:}

\paragraph{Base}~\\
\textbf{text:} Na pequena aldea de Cambados, hai un xaleiro chamado Pipo que é famoso polo seu poboalón. Cada día, moitos habitantes da aldea van a pedirlle que  [...] historia original sobre a aprendizaxe e a transmisión da arte do xaleo en Cambados, unha aldea galega. Espero que gostes desta historia en galego.\\
\textbf{ttr:} 0.5047923322683706\\
\textbf{entropy:} 6.70327761683191\\
\textbf{ngram\_overlap:} 0.0\\
\textbf{freq\_target:} 0.8466453674121406\\
\textbf{freq\_comparison:} \{\textquotesingle{}es\textquotesingle{}: 0.5846645367412141, \textquotesingle{}fr\textquotesingle{}: 0.33226837060702874\}

\paragraph{Entrenado}~\\
\textbf{text:} A historia debe ser orixinal e coherente, con personaxes, situacións e acontecementos reais, e con elementos de humor.\\
\textbf{ttr:} 0.8333333333333334\\
\textbf{entropy:} 3.794653473544342\\
\textbf{ngram\_overlap:} 0.0\\
\textbf{freq\_target:} 1.0\\
\textbf{freq\_comparison:} \{\textquotesingle{}es\textquotesingle{}: 0.7222222222222222, \textquotesingle{}fr\textquotesingle{}: 0.2777777777777778\}

\subsection{Tarea: TRADUCCION}\label{tarea-traduccion}
\begin{longtable}[]{@{}lll@{}}
\toprule\noalign{}
Métrica & Base & Entrenado \\
\midrule\noalign{}
\endhead
\bottomrule\noalign{}
\caption{Métricas de rendimiento para la tarea Traduccion (Gallego QLoRA).}
\label{tab:metricas_traduccion_Gallego QLoRA}
\endlastfoot
BLEU & 23.75 & \textbf{\underline{32.75}} \\
chrF & \textbf{\underline{51.92}} & 51.50 \\
\end{longtable}

\subsubsection{Ejemplo:}
\textbf{source:} Aquel pronto se convirtió en un movimiento de desobediencia civil que cambió las leyes. \\
\textbf{reference:} Aquela arroutada converteuse nun movemento de disobediencia civil que cambiou as leis. 

\paragraph{Base}~\\
\textbf{translated:} Ese pronto se convirxe en un movimento de desobediencia cívica que mudou as leis. A

\paragraph{Entrenado}~\\
\textbf{translated:} 

\subsection{Tarea: ROUND\_TRIP}\label{tarea-round_trip}
\begin{longtable}[]{@{}lll@{}}
\toprule\noalign{}
Métrica & Base & Entrenado \\
\midrule\noalign{}
\endhead
\bottomrule\noalign{}
\caption{Métricas de rendimiento para la tarea Round Trip (Gallego QLoRA).}
\label{tab:metricas_round_trip_Gallego QLoRA}
\endlastfoot
BLEU & 17.55 & \textbf{\underline{48.94}} \\
chrF & 41.70 & \textbf{\underline{50.43}} \\
\end{longtable}

\subsubsection{Ejemplo:}
\textbf{intermediate\_language:} es\\
\textbf{source:} Aínda que o inverno foi rigoroso, as flores xurdiron a moreas na primavera, tapizando os campos cun colorido que semellaba unha compensación da natureza. 

\paragraph{Base}~\\
\textbf{intermediate:} Aunque el invierno fue riguroso, las flores comenzaron a brotar en multitudes durante la primavera, cubriendo los campos con un colorido que parec [...] flores comenzaron a brotar en multitudes durante la primavera, cubriendo los campos con un colorido que parecía una compensación de la naturaleza.\\
\textbf{translated:} Aunque o inverno foi rígruoso, as flores comezaron a brotar en multitudes durante a primaver, cubrindo os campos co un colorido que parecía unha c [...] ruoso, as flores comezaron a brotar en multitudes durante a primaver, cubrindo os campos co un colorido que parecía unha compensación da natureza.

\paragraph{Entrenado}~\\
\textbf{intermediate:} \\
\textbf{translated:} 

\subsection{Tarea: VOCABULARIO}\label{tarea-vocabulario}
\begin{longtable}[]{@{}lll@{}}
\toprule\noalign{}
Métrica & Base & Entrenado \\
\midrule\noalign{}
\endhead
\bottomrule\noalign{}
\caption{Métricas de rendimiento para la tarea Vocabulario (Gallego QLoRA).}
\label{tab:metricas_vocabulario_Gallego QLoRA}
\endlastfoot
accuracy & \textbf{\underline{0.09}} & 0.03 \\
accuracy\_lower & \textbf{\underline{0.10}} & 0.03 \\
levenshtein & \textbf{\underline{4.93}} & 5.50 \\
\end{longtable}

\subsubsection{Ejemplo:}
\textbf{masked\_sentence:} A miña avoa {\textless}mask{\textgreater} onte, mais deixounos cheos de lembranzas fermosas.\\
\textbf{missing\_word:} espichou

\paragraph{Base}~\\
\textbf{predicted:} foi

\paragraph{Entrenado}~\\
\textbf{predicted:} 

\subsection{Tarea: ORTOGRAFIA}\label{tarea-ortografia}
\begin{longtable}[]{@{}lll@{}}
\toprule\noalign{}
Métrica & Base & Entrenado \\
\midrule\noalign{}
\endhead
\bottomrule\noalign{}
\caption{Métricas de rendimiento para la tarea Ortografia (Gallego QLoRA).}
\label{tab:metricas_ortografia_Gallego QLoRA}
\endlastfoot
BLEU & 51.70 & \textbf{\underline{61.28}} \\
chrF & 79.71 & \textbf{\underline{80.51}} \\
Levenshtein & \textbf{\underline{13.99}} & 16.99 \\
precision & 0.24 & \textbf{\underline{0.34}} \\
recall & 0.62 & \textbf{\underline{0.64}} \\
F1 & 0.28 & \textbf{\underline{0.35}} \\
errores\_totales & 2.00 & 2.00 \\
errores\_corregidos & 0.93 & \textbf{\underline{1.02}} \\
errores\_no\_corregidos & 1.07 & \textbf{\underline{0.98}} \\
errores\_nuevos & 2.88 & \textbf{\underline{1.88}} \\
\end{longtable}

\subsubsection{Ejemplo:}
\textbf{original:} A súa última palabra foi non , así que morra o conto.\\
\textbf{incorrect:} A {\textless}err t=ort{\textgreater}sua{\textless}/err{\textgreater} última palabra foi non , {\textless}err t=add{\textgreater}muy {\textless}/err{\textgreater}{\textless}err t=reord{\textgreater}morra o conto así que{\textless}/err{\textgreater}.

\paragraph{Base}~\\
\textbf{corrected:} A su última palabra foi non, moi morra o conto así que

\paragraph{Entrenado}~\\
\textbf{corrected:} A súa última palabra foi non, moito morra o conto así que